\documentclass[conference]{IEEEtran}

\usepackage{cite}
\usepackage{amsmath,amssymb,amsfonts}
\usepackage{graphicx}
\usepackage{textcomp}
\usepackage{xcolor}
\usepackage{booktabs}
\usepackage{url}


\begin{document}

\title{A Composite Spatial Accessibility Index for Urban Facility Location}

\author{\IEEEauthorblockN{Rasmus Erlemann, Hruthika Goud Tigulla, Shashank Gajula}
\IEEEauthorblockA{\textit{College of Computing and Informatics}\\
\textit{University of North Carolina at Charlotte}\\
Charlotte, North Carolina, USA\\
\{rerlemann, htigulla, sgajula3\}@charlotte.edu}}

\maketitle

% =============================================================
\begin{abstract}
% =============================================================
Choosing where to open a new store or clinic is a decision that affects
thousands of people, yet the three factors that matter most are almost
never analyzed together: how many potential customers live nearby, how
quickly they can drive there on real roads, and how many competing
facilities already serve the area.  We introduce a composite
accessibility index that combines all three into a single bounded,
weight-controllable score for every candidate location, letting
decision-makers compare sites on equal footing.

The index fuses the Huff gravity model (consumer demand capture), a
road-network reachability score (average driving time weighted by
population), and a population-normalized count of existing facilities.
Computing road travel times at city scale is the main computational
bottleneck.  We show that by exploiting the symmetry of undirected road
graphs, shortest paths can be computed from the smaller of the two
source sets in a single compiled operation, achieving 4.6$\times$ speedup
for Charlotte, NC over running one query per candidate site.

The framework is validated through four experiments: runtime scaling
across three US metropolitan road networks of increasing size (V1),
weight calibration against the known store network of an established grocery
retailer (V2), a head-to-head comparison with routing methods from three
recent published studies (V3), and a permutation significance test
confirming the calibration is non-random (V4, $p = 0.0099$).  The complete pipeline uses only publicly
available OpenStreetMap road data and US Census demographics, and runs
in standard Python without external routing infrastructure.
\end{abstract}

\begin{IEEEkeywords}
facility location, spatial accessibility, Huff gravity model,
Dijkstra shortest path, geographic information systems, urban planning
\end{IEEEkeywords}

% =============================================================
\section{Introduction}
% =============================================================

The Facility Location Problem asks which candidate site best balances
demand, travel convenience, and market competition.  Classical
formulations such as the $p$-median problem~\cite{ReVelle2005} assume
straight-line distances, but circuitous urban road networks can change
which site ranks highest.  Evaluating candidates on actual road-network
travel time therefore requires shortest-path distances between every
candidate and every residential demand zone, a computation whose cost
grows with network size and candidate count.

Gravity-based models, particularly the Huff model~\cite{Huff1964},
allocate demand probabilistically across competing facilities, but most
implementations rely on Euclidean distances and are disconnected from
network accessibility measures.  Network accessibility
studies~\cite{Hansen1959, Xu2020} quantify reachability but rarely
account for market dynamics.  Gravity models consequently overestimate
accessibility in poorly connected areas, while accessibility studies
ignore market saturation.

We introduce a site selection framework that combines Huff market
potential, a Hansen road-network accessibility index, and a
stores-per-capita competition metric into a single bounded score.
Decision-makers can compare thousands of candidates and adjust scoring
weights to match their priorities, without GIS expertise or external
routing infrastructure.  The accessibility computation uses
Multi-source Dijkstra via scipy's compiled C backend, exploiting the
fact that on undirected road graphs distances can be computed from
whichever source set is smaller---this is a standard property of
undirected graphs, applied here to make the framework tractable at city
scale (Section~\ref{sec:method}).  We do not claim Reverse Dijkstra as
a novel algorithm; our contribution is its systematic application within
a composite scoring pipeline that is validated quantitatively.

The main contributions of this paper are:
\begin{enumerate}
    \item A composite spatial accessibility index integrating Huff
    market potential, Hansen road-network accessibility, and a
    population-normalized competition penalty into a single
    weight-controllable candidate score (Section~\ref{sec:method}).
    \item A reproducible pipeline using only public OpenStreetMap
    and US Census data, running on commodity hardware without
    proprietary routing software or GIS infrastructure.
    \item A four-experiment validation protocol (V1--V4): runtime
    scaling across three US cities, retrospective weight calibration
    against a real store network, a head-to-head routing comparison
    with three published methods, and a permutation significance test
    confirming calibration is non-random ($p = 0.0099$).
    \item Empirical evidence that Multi-source Dijkstra from the
    smaller source set achieves $1.9$--$4.6\times$ speedup over
    standard sequential Dijkstra and outperforms OSRM Contraction
    Hierarchies on one-to-all accessibility workloads.
\end{enumerate}

The full pipeline runs on commodity hardware using only publicly
available Census and OpenStreetMap data, with no proprietary routing
engine, GIS license, or external infrastructure.

% =============================================================
\section{Related Work}
% =============================================================

\subsection{Facility Location and Gravity Models}

Classical facility location models ($p$-median~\cite{Hakimi1964},
$p$-centre, maximal covering~\cite{Church1974}, set
covering~\cite{Toregas1971}) provide rigorous
optimization frameworks but rely on simplified distance metrics and
homogeneous demand~\cite{ReVelle2005, Owen1998, Farahani2012,
Drezner2024}.  Equity considerations have also been incorporated into
location planning~\cite{Morrill1977}.  The Huff
model~\cite{Huff1964} treats consumer choice probabilistically,
allocating demand inversely proportional to a power of travel
distance~\cite{Wilson1971, Nakanishi1974, Fotheringham1983,
Suhara2019, Liang2020}.  Network-aware extensions account for
patronage patterns shaped by competitor distribution and transportation
structure~\cite{Sevtsuk2017, Verhetsel2022}, but most implementations
still use Euclidean impedance and do not integrate accessibility as a
separate scored component.

\subsection{Network Accessibility and Travel-Time Realism}

Hansen~\cite{Hansen1959} defined accessibility as the ease with which
opportunities can be reached, a concept now standard in transportation
planning~\cite{Geurs2004,Luo2003,Wang2012}.  Empirical work shows that
accessibility patterns follow spatial laws shaped by road network
structure~\cite{Xu2020} and that socioeconomic context exposes
inequities hidden under pure travel-time measures~\cite{Kharel2024},
motivating demand weights that capture both volume and purchasing power.

\subsection{Competition and Market Saturation in Site Selection}

Market saturation is typically omitted from accessibility models.
Food desert studies document how grocery access varies by neighbourhood
income and race~\cite{Apparicio2007,Beaulac2009,VerPloeg2012}, but
focus on measuring deficits rather than scoring candidate sites.  A
population-normalized competition index fills this gap, capturing
whether supply keeps pace with demand in each candidate's catchment.

\subsection{Shortest-Path Computation in Accessibility Studies}

Point-to-point acceleration algorithms such as A*~\cite{Hart1968},
Bidirectional Dijkstra, and Contraction Hierarchies
(CH)~\cite{Geisberger2008} exploit goal-directed pruning that provides
no benefit when the full distance vector from a source is required, as
in accessibility computation.  A comprehensive survey of route
planning methods is given in~\cite{Bast2016}.  The correct one-to-all
extension of Contraction Hierarchies is PHAST~\cite{Delling2011},
which requires specialized C++ implementations not available in
standard scientific Python toolkits~\cite{Virtanen2020}.

Recent accessibility studies have used different strategies.
Jin and Lu~\cite{Jin2022} used ArcGIS OD Cost Matrix (sequential
Dijkstra) for multi-mode food accessibility analysis.  Kang et
al.~\cite{Kang2020} parallelized the Enhanced Two-Step Floating
Catchment Area method across CPU cores.  Horton et
al.~\cite{Horton2025} used OSRM with Contraction Hierarchies for a
Kolm-Pollak equity optimization study.  We show that an efficient
implementation choice---exploiting undirected-graph symmetry via
Multi-source Dijkstra in standard Python---achieves comparable or
better performance than external routing engines, without preprocessing
infrastructure or proprietary software.

% =============================================================
\section{Methodology}
\label{sec:method}
% =============================================================

Fig.~\ref{fig:pipeline} summarises the full system pipeline.  Raw
inputs from three public data sources feed into preprocessing stages
(Extract, Transform, Load (ETL), spatial standardization, road graph construction, demand
attribute normalization, and demand-to-graph mapping), which supply
three parallel scoring components (Huff market potential, Multi-source
Dijkstra accessibility, and local competition index).  The components
are combined into a final weighted suitability score and filtered by a
diversity constraint to produce the ranked location recommendations.

\begin{figure*}[!t]
    \centering
    \includegraphics[width=\textwidth]{archimate_exact_v2.png}
    \caption{ArchiMate~3.x layered architecture of the proposed framework.
    \textbf{Business layer}: Retail Strategist $\to$ Site Selection Report
    $\to$ Top-$N$ Locations.
    \textbf{Application layer}: ETL $\to$ Road Graph Construction $\to$
    Multi-source Dijkstra produces a $559\times282{,}751$ Distance Matrix;
    three independent scoring components (Huff Market Capture, Accessibility
    Index, Competition Index) feed Composite Scoring.
    \textbf{Technology layer}: HPC node (Intel~i9, 80\,GB,
    Python/SciPy) with CSR sparse matrix storage.
    Dashed arrows: data access; open arrows: serving; filled-circle arrows:
    assignment.}
    \label{fig:pipeline}
\end{figure*}


\subsection{Data Sources and Preprocessing}

The framework requires three publicly available inputs:

\textbf{Demand zones.}  Census block groups (BGs) from the American
Community Survey~\cite{Census2023}, each carrying population
($\text{pop}_i$), median household income ($\text{inc}_i$), and
population density ($\text{dens}_i$).  For Charlotte, NC (Mecklenburg County):
$N_{\text{BG}} = 559$ block groups after clipping to the city
boundary.  Fig.~\ref{fig:block_groups} shows the 559 block group
boundaries across Mecklenburg County.

\begin{figure}[!t]
    \centering
    \includegraphics[width=\linewidth]{block_group.png}
    \caption{Census block group boundaries across Mecklenburg County
    (559 groups after clipping to the city boundary).  Each polygon
    carries four ACS attributes: Geographic Identifier (GEOID), total population, area
    (sq.\ mi), population density (persons/sq.\ mi), and median
    household income.  These attributes feed directly into the
    composite demand weight $B_i$ (Eq.~\ref{eq:demand_weight}).}
    \label{fig:block_groups}
\end{figure}

\textbf{Road network.}  OpenStreetMap~\cite{OSM2024} drive-accessible
road segments, downloaded using OSMnx~\cite{Boeing2017} and
represented as an undirected weighted graph $G = (V, E, w)$ where edge
weight $w_e$ is segment length in metres.  Charlotte:
$|V| = 282{,}751$ nodes, $|E| = 585{,}666$ edges.  The graph is
stored as a symmetric Compressed Sparse Row (CSR) matrix and passed
directly to \texttt{scipy.sparse.csgraph.dijkstra}~\cite{Virtanen2020}
with \texttt{directed=False}.

\textbf{Candidate sites.}  Extracted from OSM commercial landuse
polygons (\texttt{landuse=retail/commercial},
\texttt{shop=supermarket}), area-filtered to
$500$--$500{,}000\,\text{m}^2$ (lower bound excludes kiosks and
parking stalls; upper bound excludes large industrial parcels
unsuitable for retail), represented by polygon centroids.
Charlotte: $N_{\text{cands}} = 2{,}450$.  All 28 Harris Teeter
ground-truth stores fall within 0.5 miles of a candidate.

\textbf{Existing facilities.}  Coordinates of target-brand and
competitor stores.  Charlotte: 28 Harris Teeter plus 52 competitors
(Kroger, Publix, Food Lion, Aldi, Lidl, and similar full-service
supermarkets), totaling 80 grocery stores within Mecklenburg County.
Competitor locations were extracted from OSM nodes tagged
\texttt{shop=supermarket} or \texttt{shop=grocery} within the county
boundary and verified against publicly available store locator data.

All geometries are projected to NAD83/North Carolina (EPSG:32119) for
metric computation.

\textbf{Node snapping.}  Block group centroids and candidate sites are snapped to their nearest road node via k-d tree ($k=1$), producing index arrays \texttt{bg\_node\_ids} and \texttt{cand\_node\_ids}.  Block groups with suppressed ACS income values are excluded from downstream computation.

\begin{table}[!t]
\centering
\caption{Charlotte Case Study: Data Summary}
\label{tab:data_summary}
\begin{tabular}{lr}
\toprule
\textbf{Component} & \textbf{Count} \\
\midrule
Census block groups (demand zones) & 559 \\
Candidate locations (OSM commercial) & 2{,}450 \\
Harris Teeter stores (ground truth) & 28 \\
Competitor grocery stores & 52 \\
Road network nodes & 282{,}751 \\
Road network edges & 585{,}666 \\
\bottomrule
\end{tabular}
\end{table}

\subsection{Normalization}
\label{sec:norm}

Before any component is used in scoring, each raw variable is
independently scaled to the interval $[0.01, 1.0]$.  Raw values are
first clipped to their 1st and 99th percentiles to suppress the effect
of extreme outliers, then mapped by a linear min-max transform:
\begin{equation}
\tilde{x}_i
  = 0.01 + 0.99\,\frac{
      \text{clip}(x_i,\,p_1,\,p_{99}) - p_1
    }{
      p_{99} - p_1
    }
\label{eq:norm}
\end{equation}
where $p_1$ and $p_{99}$ are the 1st and 99th percentiles of variable
$x$.  The lower bound of 0.01 (rather than 0) ensures that no element
receives a zero weight, which would remove it entirely from
multiplicative terms.  This normalization is applied independently to
$\text{pop}_i$, $\text{inc}_i$, $\text{dens}_i$, and to each of the
three scoring components $P_k$, $A_k$, $S_k$ before combination in
Eq.~(\ref{eq:score}).  The tilde notation $\tilde{(\cdot)}$ denotes a
value that has been normalized by Eq.~(\ref{eq:norm}).

\subsection{Demand Weight Construction}

Each block group $i$ receives a composite demand weight $B_i$
constructed from three ACS attributes.  Let $\widetilde{\text{pop}}_i$
denote the normalized total population (volume of potential customers),
$\widetilde{\text{inc}}_i$ the normalized median household income
(purchasing power), and $\widetilde{\text{dens}}_i$ the normalized
population density in residents per square mile (spatial concentration
of demand), each scaled by Eq.~(\ref{eq:norm}).  Then:
\begin{equation}
B_i = 0.4\,\widetilde{\text{pop}}_i
    + 0.3\,\widetilde{\text{inc}}_i
    + 0.3\,\widetilde{\text{dens}}_i
\label{eq:demand_weight}
\end{equation}
The default coefficients $(0.4,\;0.3,\;0.3)$ are configurable inputs,
not empirically derived constants.  A discount chain entering a new
market might increase the population coefficient because volume
outweighs affluence; a premium grocer might increase the income
coefficient because purchasing power drives basket size.

\subsection{Huff Market Potential}

Let $W_1$, $W_2$, $W_3 \geq 0$ with $W_1 + W_2 + W_3 = 1$ denote the
outer scoring weights for demand, accessibility, and competition
respectively; these are defined formally in the composite score
(Eq.~\ref{eq:score}) and calibrated retrospectively in
Section~\ref{sec:v2}.

The Huff gravity model~\cite{Huff1964} estimates the probability that
demand from block group $i$ is captured by candidate $k$ rather than
by existing competitors $j \in C$:
\begin{equation}
\text{Share}_{ik} = \frac{t_{ik}^{-\beta}}
                         {t_{ik}^{-\beta}
                          + \sum_{j \in C_K(i)} t_{ij}^{-\beta}}
\label{eq:huff}
\end{equation}
where $t_{ik}$ is travel time in minutes (Euclidean distance divided by
an assumed speed of 35\,mph), $\beta$ is the distance-decay exponent
(default 2.0), and $C_K(i)$ is the set of $K$ nearest competitor stores
to block group $i$ (default $K = 3$).  Using $K$ nearest competitors
rather than all competitors avoids diluting the denominator with distant
facilities that have negligible influence on local choice behavior.

The raw market potential for candidate $k$ is:
\begin{equation}
P_k = \sum_{i \in R_k} B_i\,\text{Share}_{ik}
\label{eq:potential}
\end{equation}
where $R_k$ is the set of block groups whose centroids fall within a
5-mile catchment radius of candidate $k$, and $B_i$ is the composite
demand weight from Eq.~(\ref{eq:demand_weight}).

The Huff model uses Euclidean distance rather than road-network
distance.  This choice is deliberate: Huff captures consumer
\textit{perception} of effort, not actual routing cost.  Shoppers
comparing two stores reason in straight-line terms, not turn-by-turn
directions.  Huff's original formulation used Euclidean distance for
this reason~\cite{Huff1964}, and subsequent retail gravity
implementations have largely followed suit~\cite{Suhara2019,
Liang2020}.  Road-network distance is the appropriate metric for the
accessibility component (Eq.~\ref{eq:tbar}), which quantifies travel
burden on real infrastructure.

\subsection{Road-Network Accessibility Index}
\label{sec:sssp}

\subsubsection{Computing the Distance Matrix}

The accessibility component requires one-to-all shortest-path
distances from each source to all road nodes.  Three strategies are
possible:

\textbf{Standard (Forward) Dijkstra.}  Run one shortest-path
computation from each of the $N_{\text{cands}}$ candidates, requiring
$N_{\text{cands}}$ calls.

\textbf{Reverse Dijkstra.}  Because $G$ is undirected, $d(A,B) =
d(B,A)$.  Running from each of the $N_{\text{BG}}$ block groups
instead yields the same distance matrix with fewer calls when
$N_{\text{BG}} < N_{\text{cands}}$, giving a theoretical call-count
ratio of $N_{\text{cands}}/N_{\text{BG}} = 2{,}450/559 \approx 4.38\times$
for Charlotte.

\textbf{Multi-source Dijkstra.}  Rather than issuing one call per
block group, all $N_{\text{BG}}$ source nodes (stored in
\texttt{bg\_node\_ids}) are passed simultaneously to a single
\texttt{scipy.sparse.csgraph.dijkstra}~\cite{Virtanen2020} call with
\texttt{directed=False} and \texttt{indices=bg\_node\_ids}.  The
function returns a dense matrix $D$ of shape $N_{\text{BG}} \times
|V|$, where $D[s, v]$ holds the shortest-path distance from source
block group $s$ to road node $v$.

This is mathematically equivalent to $N_{\text{BG}}$ sequential reverse-Dijkstra calls and faster because (a) $N_{\text{BG}} < N_{\text{cands}}$ for Charlotte and (b) a single compiled C call eliminates Python interpreter overhead across hundreds of invocations.

The road distance from block group $i$ to candidate $k$ is then
retrieved as a single array lookup:
\begin{equation}
d_{ik}^{\text{road}} = D\bigl[\,\text{bg\_idx}(i),\;
                                \text{cand\_node}(k)\,\bigr]
\label{eq:lookup}
\end{equation}
and converted to travel time in minutes:
\begin{equation}
t_{ik}^{\text{road}}
  = \frac{d_{ik}^{\text{road}}}{35 \times 0.44704 \times 60}
  = \frac{d_{ik}^{\text{road}}}{936.8}\;\text{[minutes]}
\label{eq:ttime}
\end{equation}
where $d_{ik}^{\text{road}}$ is in metres, $35\,\text{mph} \times
0.44704\,\text{m\,s}^{-1}\text{mph}^{-1} = 15.65\,\text{m/s}$, and
dividing by 60 converts seconds to minutes.

The Dijkstra distance limit is set to $r \times 1.30$ (circuity
factor) to accommodate road-network indirection beyond the straight-
line catchment radius $r = 5$ miles.

\subsubsection{Accessibility Score Computation}

The demand-weighted average shortest-path travel time from all block
groups in the catchment to candidate $k$ is:
\begin{equation}
\bar{t}_k = \frac{\sum_{i \in R_k} B_i \cdot t_{ik}^{\text{road}}}
                  {\sum_{i \in R_k} B_i}
\label{eq:tbar}
\end{equation}
where $t_{ik}^{\text{road}}$ is the Dijkstra shortest-path travel time
on $G$ and $B_i$ is the composite demand weight from
Eq.~(\ref{eq:demand_weight}).  Block groups with no finite path to
candidate $k$ are excluded from both the numerator and denominator.

The accessibility score follows the Hansen~\cite{Hansen1959}
negative-exponential form:
\begin{equation}
A_k = \exp\!\bigl(-\alpha \cdot \bar{t}_k\bigr), \quad \alpha = 0.05
\label{eq:access}
\end{equation}
The value $\alpha = 0.05$ is consistent with driving-mode decay rates
reported in the urban accessibility literature for catchment radii of
5--10\,miles~\cite{Luo2003, Wang2012}, and produces meaningful spread
across typical urban travel times: $A_k = 0.78$ at $\bar{t}_k =
5$\,min, 0.61 at 10\,min, 0.47 at 15\,min, and 0.37 at 20\,min.
Sensitivity analysis over $\alpha \in \{0.03,\,0.05,\,0.08\}$
showed Spearman rank correlation above 0.96 across all candidate
rankings, confirming the final ranking is stable with respect to
this parameter.

\subsection{Competition Saturation Index}

Within a 5-mile radius of candidate $k$, let $C_k$ be the count of all
existing stores (target-brand and competitors, since the market is
saturated regardless of brand) and $\text{Pop}_k$ the raw total
population across all block groups in that catchment:
\begin{equation}
S_k = \frac{C_k}{\text{Pop}_k / 10{,}000}
\label{eq:competition}
\end{equation}
Higher values indicate more saturated markets.  Normalizing by
population rather than area avoids penalizing dense urban sites simply
because they have many nearby stores; what matters is whether supply
keeps pace with demand.

\subsection{Composite Suitability Score}

All three components are independently normalized to $[0.01, 1.0]$ via
Eq.~(\ref{eq:norm}) before combination.  The scoring weights $W_1$,
$W_2$, $W_3 \geq 0$ with $W_1 + W_2 + W_3 = 1$ control each
component's contribution:
\begin{equation}
\text{Score}_k
  = W_1\,\widetilde{P}_k
  + W_2\,\widetilde{A}_k
  - W_3\,\widetilde{S}_k
\label{eq:score}
\end{equation}
where $\widetilde{P}_k$, $\widetilde{A}_k$, and $\widetilde{S}_k$ are
the normalized market potential, accessibility score, and competition
index for candidate $k$, each produced by Eq.~(\ref{eq:norm}).  Market
potential and accessibility contribute positively; competition enters
as a penalty.  The score range is $\text{Score}_k \in [-W_3,\;
W_1 + W_2]$.

The weights $W_1, W_2, W_3$ are calibrated against the known store
network of whichever facility type is being analyzed
(Section~\ref{sec:v2}), making the framework adaptable to any retailer,
healthcare provider, or public service.

\subsection{Spatial Diversity Selection}

The top $N$ sites are selected by greedy spatial-diversity search: the
highest-scoring candidate is chosen first, then iteratively the next
highest-scoring candidate at least $d_{\min} = 2.5$ miles from every
already-selected site.  This heuristic prevents the clustering that
unconstrained score-maximization would produce, spreading
recommendations across the market area.

\begin{table}[!t]
\centering
\caption{Model Parameters and Default Values}
\label{tab:params}
\begin{tabular}{lll}
\toprule
\textbf{Symbol} & \textbf{Default} & \textbf{Description} \\
\midrule
$\beta$    & 2.0       & Huff distance-decay exponent \\
$\alpha$   & 0.05      & Hansen accessibility decay rate \\
$K$        & 3         & Nearest competitors in Huff denominator \\
$r$        & 5 miles   & Catchment radius \\
$d_{\min}$ & 2.5 miles & Min.\ separation for diversity \\
Speed      & 35 mph    & Assumed uniform driving speed \\
Circuity   & 1.30      & Dijkstra limit $= r \times$ circuity \\
$W_1, W_2, W_3$ & 0.4,\;0.3,\;0.3 & Component weights (sum to 1) \\
\bottomrule
\end{tabular}
\end{table}

% =============================================================
\section{Case Study: Charlotte, NC}
\label{sec:results}
% =============================================================

We applied the framework to the Charlotte, NC metropolitan region using
the data summarised in Table~\ref{tab:data_summary}.

\subsection{Demographic Demand Structure}

Population density concentrates in central Charlotte and southeastern
corridors; median household income does not track density, with
high-density inner-city areas often coinciding with lower incomes.
Table~\ref{tab:demo} reports the demographic spread.

\begin{table}[!t]
\centering
\caption{Demographic Profile of Charlotte Block Groups (559 BGs)}
\label{tab:demo}
\begin{tabular}{lrrr}
\toprule
\textbf{Attribute} & \textbf{Median} & \textbf{10th pct.}
                   & \textbf{90th pct.} \\
\midrule
Population                  & 1{,}412 & 612    & 2{,}891    \\
Median income (USD)         & 62{,}400 & 28{,}100 & 113{,}500 \\
Density (per sq.\ mi)       & 3{,}210 & 640    & 9{,}870    \\
\bottomrule
\end{tabular}
\end{table}

Weighting by population alone would misrepresent market quality.  The
composite weight $B_i$ from Eq.~(\ref{eq:demand_weight}) captures both
volume and economic capacity simultaneously.

\subsection{Network Accessibility Patterns}

Accessibility scores from Eq.~(\ref{eq:access}) depend on network
topology.  Candidates near arterials and highway interchanges
(I-77, I-85, I-485) score highest; 61.5\% of the 2{,}450 candidates
fall in the 8--16\,min travel-time band ($A_k \in [0.45, 0.67]$),
while only 8.8\% exceed 16\,min ($A_k < 0.45$).  Some geographically
central locations score poorly due to circuitous local streets, while
peripheral sites near highway access points outperform them.

\subsection{Competitive Saturation}

Established commercial corridors in south and northeast Charlotte
exhibit competition indices $S_k > 1.5$ stores per 10{,}000 people.
Without the $W_3$ penalty, these saturated areas would dominate
rankings purely due to demand and accessibility.  The competition
component redirects recommendations toward emerging suburban zones
($S_k < 0.3$) with growing populations and limited existing facilities.

\subsection{Component Independence}
\label{sec:component_analysis}

A three-component framework is justified only if the components capture
different spatial information.  Table~\ref{tab:component_stats} reports
descriptive statistics for each raw component across all 2{,}450
candidates, and Table~\ref{tab:correlation} reports pairwise Spearman
rank correlations after normalization.

\begin{table}[!t]
\centering
\caption{Raw Component Statistics Across 2,450 Candidates}
\label{tab:component_stats}
\begin{tabular}{lrrrr}
\toprule
\textbf{Component} & \textbf{Mean} & \textbf{Std} & \textbf{Range}
                    & \textbf{CV} \\
\midrule
$P_k$ (Market Potential) & 4.394 & 1.834 & 7.712 & 0.417 \\
$A_k$ (Accessibility)    & 0.709 & 0.022 & 0.137 & 0.030 \\
$S_k$ (Competition)      & 0.848 & 0.158 & 0.869 & 0.186 \\
\bottomrule
\end{tabular}
\end{table}

Market potential shows the strongest differentiation (CV $= 0.417$); competition moderate (CV $= 0.186$); raw accessibility narrow (CV $= 0.030$), reflecting Charlotte's uniform network.  Percentile-robust normalization maps all three into $[0.01, 1.0]$, ensuring each contributes meaningful spread to the final score.

\begin{table}[!t]
\centering
\caption{Pairwise Spearman Rank Correlation (Normalized Components)}
\label{tab:correlation}
\begin{tabular}{lccc}
\toprule
& $\widetilde{P}$ & $\widetilde{A}$ & $\widetilde{S}$ \\
\midrule
$\widetilde{P}$ (Market Potential) & 1.000 & 0.383 & 0.065 \\
$\widetilde{A}$ (Accessibility)    & 0.383 & 1.000 & 0.040 \\
$\widetilde{S}$ (Competition)      & 0.065 & 0.040 & 1.000 \\
\bottomrule
\end{tabular}
\end{table}

All pairwise correlations fall below 0.4: competition is nearly orthogonal to both demand and accessibility ($\rho < 0.07$), and the moderate $\widetilde{P}$--$\widetilde{A}$ correlation ($\rho = 0.383$) is far from deterministic.  Variance decomposition under balanced weights confirms each component contributes independently:

\begin{table}[!t]
\centering
\caption{Score Variance Decomposition
         ($W_1 = 0.4$, $W_2 = 0.3$, $W_3 = 0.3$)}
\label{tab:variance}
\begin{tabular}{lrr}
\toprule
\textbf{Term} & \textbf{Variance}
              & \textbf{\% of $\text{Var}(\text{Score})$} \\
\midrule
$W_1 \widetilde{P}$ (Market Potential) & 0.01164 & 44.5\% \\
$W_2 \widetilde{A}$ (Accessibility)    & 0.00435 & 16.6\% \\
$W_3 \widetilde{S}$ (Competition)      & 0.00476 & 18.2\% \\
Cross-covariance terms                 & 0.00537 & 20.6\% \\
\midrule
$\text{Var}(\text{Score})$             & 0.02612 & 100.0\% \\
\bottomrule
\end{tabular}
\end{table}

Market potential explains 44.5\%, competition 18.2\%, and accessibility 16.6\% of score variance; none is negligible.  The remaining 20.6\% arises from cross-covariance between components (primarily the moderate $\widetilde{P}$--$\widetilde{A}$ correlation, $\rho = 0.383$), meaning the components are not fully independent but each retains unique explanatory content.

\subsection{Composite Suitability and Top Locations}

\begin{figure}[!t]
    \centering
    \includegraphics[width=\linewidth]{top_locations.png}
    \caption{Top-10 recommended facility locations across Charlotte
    under balanced weights ($W_1 = 0.4$, $W_2 = 0.3$, $W_3 = 0.3$).
    Large blue circles mark the recommended sites; surrounding rings
    indicate their 5-mile catchment search zones.  Small teal markers
    show existing Harris Teeter stores and competitor grocery stores.
    Recommended sites concentrate in north and south Charlotte near the
    I-485 corridor, in areas combining strong demand, good road
    accessibility, and limited existing grocery coverage.}
    \label{fig:top_locations}
\end{figure}

High-scoring regions form spatial clusters where demand, accessibility,
and low competition coincide.  Under balanced weights
($W_1 = 0.4$, $W_2 = 0.3$, $W_3 = 0.3$), top-ranked sites combine
market potential above 0.75, accessibility above 0.70, and competition
below 0.40 stores per 10{,}000 people.  Fig.~\ref{fig:top_locations}
shows the top-10 locations; recommended sites concentrate in north and
south Charlotte near the I-485 corridor, in areas with strong demand
and limited existing coverage.



% =============================================================
\section{Validation}
\label{sec:validation}
% =============================================================

Four experiments validate the framework.  V1 benchmarks multi-city runtime scaling.  V2 calibrates weights against Harris Teeter's known store network.  V3 compares routing methods from three published studies.  V4 tests statistical significance of the calibration.  V1 ran on an HPC node (Intel~i9, 80\,GB, RHEL~9.7, Python~3.9.23, SciPy~1.13.1); V3 ran on a workstation (Core~i5-1035G1, 8\,GB, Windows~11, Python~3.12.5, Docker~29.0.1) with all four methods sharing identical hardware; all timings are single-run measurements with caches cleared.

% ------ V1 ------
\subsection{V1: Multi-City Runtime Scaling}
\label{sec:v1}

The central computational task in our framework is building the block
group distance matrix: a dense $N_{\text{BG}} \times |V|$ matrix
containing the shortest-path distance from every block group centroid
to every road node.  Multi-source Dijkstra constructs this matrix in a
single compiled C call.  The standard baseline runs one Dijkstra call
per candidate site from the larger source set.

We measured both approaches on three US cities of increasing network
size (Table~\ref{tab:v1_scaling}).

\begin{table}[!t]
\centering
\caption{V1: Multi-City Runtime Scaling \& Speedup (HPC, 80\,GB RAM)}
\label{tab:v1_scaling}
\resizebox{\linewidth}{!}{%
\begin{tabular}{lrrrrrr}
\toprule
\textbf{City} & \textbf{Nodes} & \textbf{BGs} & \textbf{Cands}
              & \textbf{Ours} & \textbf{Standard} & \textbf{Speedup} \\
\midrule
Charlotte   & 282{,}751     & 559     & 2{,}450  & 42.7\,s    & 196.8\,s    & 4.63$\times$ \\
Atlanta     & 916{,}452     & 1{,}800 & 1{,}148  & 2.9\,min   & 5.5\,min    & 1.89$\times$ \\
Los Angeles & 1{,}018{,}083 & 6{,}400 & 14{,}208 & 37.8\,min  & 93.7\,min\textsuperscript{a} & 2.48$\times$ \\
\bottomrule
\end{tabular}}
\vspace{2pt}
\footnotesize\textsuperscript{a}LA Standard Dijkstra extrapolated from 70\% completion (3{,}934\,s elapsed, stable per-source rate).
\end{table}

Multi-source Dijkstra always runs from the smaller source set; for
Atlanta, candidates (1{,}148) are fewer than block groups (1{,}800), so
the algorithm runs from candidate nodes and the resulting distance
matrix is transposed for scoring.  Measured speedups exceed the
call-count ratio in all three cities because a single compiled C
invocation eliminates Python interpreter overhead across hundreds of
calls: Charlotte ($4.63\times$ vs.\ $4.38\times$ call-count ratio) and
Los Angeles ($2.48\times$ vs.\ $2.22\times$) benefit most; Atlanta
($1.89\times$) shows the expected result when source sets are similar in
size.

% ------ V2 ------
\subsection{V2: Retrospective Weight Calibration Against Known Store Locations}
\label{sec:v2}

V2 is a \emph{retrospective} analysis: it asks which scoring weights
best explain where Harris Teeter \emph{already} placed its 28 stores,
thereby revealing the retailer's historical siting priorities.  The
resulting calibrated weights are not used for new location
recommendations---those use the balanced default weights
($W_1=0.4$, $W_2=0.3$, $W_3=0.3$) in Section~\ref{sec:results}.

The composite score weights $W_1, W_2, W_3$ and the Huff decay
parameter $\beta$ were estimated by minimizing the mean rank of 28
known Harris Teeter store locations within the scored candidate pool
of 2{,}450 sites.  The optimization used Differential
Evolution~\cite{Storn1997} (global search, 200 iterations, population
size 12) followed by Nelder-Mead refinement~\cite{Nelder1965},
subject to $W_1 + W_2 + W_3 = 1$ and $W_i \geq 0$.

All travel times in the calibration loop were computed from the
precomputed block group distance matrix, not from Euclidean
approximations.  The calibrated weights therefore reflect actual
road-network travel patterns.

\begin{table}[!t]
\centering
\caption{V2: Calibrated Weights \& Configuration Comparison (28 HT stores, 2,450 candidates)}
\label{tab:v2_weights}
\resizebox{\linewidth}{!}{%
\begin{tabular}{lrrrrrr}
\toprule
\textbf{Configuration} & \textbf{$W_1$} & \textbf{$W_2$} & \textbf{$W_3$} & \textbf{$\beta$}
                       & \textbf{Mean Rank} & \textbf{Median Rank} \\
\midrule
Calibrated     & 0.984 & 0.016 & 0.000 & 3.0 & 921     & 712     \\
Paper defaults & 0.400 & 0.300 & 0.300 & 2.0 & 1{,}291 & 1{,}458 \\
Equal weights  & 0.333 & 0.333 & 0.333 & 2.0 & 1{,}331 & 1{,}500 \\
\bottomrule
\end{tabular}}
\end{table}

$W_1 = 0.984$ shows that Harris Teeter's siting was demand-driven: stores placed to maximise exposure to high-population, high-income areas.  The near-zero $W_3$ is consistent with a premium co-location strategy.  The calibrated $\beta = 3.0$ (vs.\ default 2.0) implies steeper distance-decay: customers prefer nearby stores even when alternatives exist.  The precomputed distance matrix reduces each calibration-loop evaluation to a column lookup ($\approx$0.4\,s per call), making the full optimization feasible in under 5 minutes.  Statistical significance of the calibration result is assessed in V4 (Section~\ref{sec:v4}).

% ------ V3 ------
\subsection{V3: Comparison with Published Routing Methods}
\label{sec:v3}

We replicated three published routing strategies on the same Charlotte network and candidate set.  Jin and Lu~\cite{Jin2022} used ArcGIS OD Cost Matrix (sequential Dijkstra: 2{,}450 \texttt{scipy} calls).  Kang et al.~\cite{Kang2020} parallelized Dijkstra across 4 cores (\texttt{multiprocessing.Pool}).  Horton et al.~\cite{Horton2025} used OSRM Contraction Hierarchies~\cite{Luxen2011} (\texttt{osrm-backend:latest}, \texttt{car.lua} profile, full North Carolina extract, 13.7M nodes).  All four methods ran on the same workstation (Core~i5-1035G1, 8\,GB, Windows~11) to ensure a fair timing comparison; OSRM requires Docker and was not available on the HPC cluster used for V1.

\begin{table}[!t]
\centering
\caption{V3: Routing Method Comparison (Charlotte, 559 BGs, 2,450 cands, 282,751 nodes)}
\label{tab:v3_comparison}
\small
\resizebox{\columnwidth}{!}{
\begin{tabular}{@{}llrrr r r@{}}
\toprule
\textbf{Method} & \textbf{Reference} & \textbf{Preproc.} & \textbf{Query} & \textbf{Total} & \textbf{Speedup} & \textbf{Spearman $\rho$} \\
\midrule
Standard Dijkstra            & Jin \& Lu 2022      & --          & 255.0\,s             & 255.0\,s  & 1.00$\times$ & baseline \\
Parallel Dijkstra (4 cores)  & Kang et al. 2020    & --          & 126.7\,s             & 126.7\,s  & 2.01$\times$ & 1.000 \\
Multi-source Dijkstra (ours) & This study          & 52.3\,s     & 0.6\,s               & 53.0\,s   & 4.81$\times$ & 1.000 \\
OSRM CH (C++)                & Horton et al. 2025  & 2{,}546\,s\textsuperscript{a} & 138.7\,s & 2{,}685\,s & 0.09$\times$ & 0.745 \\
\bottomrule
\end{tabular}
}
\vspace{2pt}
\footnotesize\textsuperscript{a}OSRM preprocessing on full NC extract (13.7M nodes) via Docker.  Query time alone (138.7\,s) exceeds our complete pipeline (53.0\,s).  Spearman $\rho$ is vs.\ Standard Dijkstra rankings.
\end{table}

Multi-source Dijkstra completes the full pipeline in 53.0\,s, the fastest of the four methods.  All Dijkstra variants produce identical rankings ($\rho = 1.000$), confirming implementation correctness.  OSRM query time alone (138.7\,s) exceeds our complete pipeline; its CH pruning provides no benefit for one-to-all accessibility workloads.  The OSRM--ours Spearman divergence ($\rho = 0.745$) reflects speed-based vs.\ distance-based edge weights---a modelling choice that affects which sites are recommended.

% ------ V4 ------
\subsection{V4: Statistical Significance of Weight Calibration}
\label{sec:v4}

To test whether the calibrated weights recover real siting preferences
beyond random chance, we ran a one-tailed permutation test:
10{,}000 times we drew 28 candidates uniformly at random from the full
pool of 2{,}450, computed their mean percentile rank (from the top,
i.e.\ higher = better ranked) under the calibrated weights
($W_1=0.984$, $W_2=0.016$, $W_3=0$, $\beta=3.0$), and compared against
the observed mean for the 28 real Harris Teeter stores.  The null
hypothesis is that real HT stores are indistinguishable from a random
draw; the alternative is that they rank higher than chance.

\begin{table}[!t]
\centering
\caption{V4: Permutation Test Results (10{,}000 draws, 28 stores, 2{,}450 candidates)}
\label{tab:v4_permutation}
\begin{tabular}{lr}
\toprule
\textbf{Statistic} & \textbf{Value} \\
\midrule
Observed mean HT percentile   & 62.9th \\
Permutation mean (null)        & 50.0th \\
Permutation std                & 5.5 \\
Draws $\geq$ observed          & 99 / 10{,}000 \\
$p$-value                      & 0.0099 \\
Significant ($\alpha = 0.05$)  & Yes \\
\bottomrule
\end{tabular}
\end{table}

The observed result (62.9th percentile from the top, meaning real HT
stores rank better than 62.9\% of all candidates) exceeds the
null-distribution mean of 50.0th with a one-tailed $p = 0.0099$: only
99 of 10{,}000 random draws matched or exceeded the actual HT store
ranks.  The calibrated framework is therefore statistically correlated
with real siting decisions at the 1\% significance level.  Note that
rank 921 of 2{,}450 in Table~\ref{tab:v2_weights} corresponds to the
62.4th percentile from the top; the 62.9th figure in Table~\ref{tab:v4_permutation}
is the average over all 28 individual store percentiles.

% =============================================================
\section{Discussion}
\label{sec:discussion}
% =============================================================

\subsection{Framework Design}

The framework scores each candidate on demand, road-network accessibility, and market competition simultaneously, combining them into a single ranked list.  Independent normalization ensures weights genuinely control each component's contribution.  Component maps can be inspected individually before combination, making recommendations auditable.

\subsection{Algorithmic Implications}

V1 and V3 confirm the same finding from different angles: running Multi-source Dijkstra from the smaller source set achieves $3$--$5\times$ speedup in US grocery retail scenarios and outperforms all four published alternatives on this one-to-all workload.  The ranking divergence between our distance-based method and OSRM's speed-based model ($\rho = 0.745$) shows that the routing backend is itself a modelling decision that practitioners should choose deliberately.

\subsection{Scalability}

The V1 benchmark confirmed the framework runs on networks from 282{,}751 nodes (Charlotte) to 1{,}018{,}083 nodes (Los Angeles), with Multi-source Dijkstra completing in under 38 minutes for the largest tested network.  Speedup varies with the ratio of block groups to candidates, not monotonically with network size.  All inputs---OSM road networks and Census ACS block groups---are publicly available for every US metropolitan area.

\subsection{Transferability}

For a new city, replacing the boundary polygon is sufficient; OSM and Census data cover every US metro area.  For a different facility type, only the facility dataset and weight calibration need to change.

\subsection{Limitations}

The current implementation assumes static demographics and a uniform
driving speed on every road edge.  The Huff model uses Euclidean
distance as a proxy for consumer perceived effort; replacing it with
road-network distance is a straightforward extension that may improve
accuracy in areas with irregular street connectivity.  The Huff model
also assumes rational distance-based choice; real consumers are
influenced by brand loyalty, product assortment, and pricing.
Catchment membership uses a binary centroid-in-radius test, so a block
group either falls entirely inside or outside the catchment.  Finally,
the greedy diversity selection is a heuristic for the NP-hard
$p$-dispersion problem and is not globally optimal.

\subsection{Future Work}

Replacing the fixed 35\,mph speed with time-of-day congestion profiles and adding a GTFS transit layer would extend the framework to multi-modal accessibility.  Replacing centroid-in-radius catchment assignment with polygon-overlap weighting would remove the artificial cliff-edge at the boundary.  Multi-city calibration could test whether the demand-first pattern seen for Harris Teeter generalises beyond Charlotte's relatively uniform road network.

% =============================================================
\section{Conclusion}
% =============================================================

The framework scores every candidate site on demand, road-network accessibility, and market competition, combining them into a single ranked list.  Four experiments test whether it works:
\begin{enumerate}
\setlength{\topsep}{0pt}
\setlength{\partopsep}{0pt}
\setlength{\itemsep}{2pt}
    \item \textbf{V1 (Runtime Scaling):}  Multi-source Dijkstra achieves
    $1.9$--$4.6\times$ speedup over standard Dijkstra across Charlotte
    (283K nodes), Atlanta (916K), and Los Angeles (1.02M).

    \item \textbf{V2 (Weight Calibration):}  Calibration against 28
    Harris Teeter stores yields $W_1 = 0.984$, $\beta = 3.0$,
    recovering the retailer's demand-first siting priorities.

    \item \textbf{V3 (Routing Comparison):}  Our method completes the
    Charlotte benchmark in 53.0\,s, faster than sequential Dijkstra
    (255.0\,s), parallel Dijkstra (126.7\,s), and OSRM-CH (2{,}685\,s).
    The OSRM divergence ($\rho = 0.745$) shows that the routing backend
    is itself a modelling decision.

    \item \textbf{V4 (Permutation Test):}  A 10{,}000-draw permutation
    test confirms the calibration is non-random ($p = 0.0099$): real
    HT stores rank at the 62.9th percentile, vs.\ 50.0th under the null.
\end{enumerate}

Variance decomposition shows all three scoring components contribute
to candidate differentiation (44.5\%, 18.2\%, 16.6\% of score
variance), with pairwise rank correlations below 0.4.

Everything runs in standard Python using only publicly available
data---OpenStreetMap for roads, US Census ACS for demographics.  No
proprietary routing engine, no GIS license, no external
infrastructure.  Pointed at any US city, the pipeline produces a
ranked candidate list without modification to the underlying code.

% =============================================================
\begin{thebibliography}{99}

\bibitem{ReVelle2005}
C.~S. ReVelle and H.~A. Eiselt, ``Location analysis: A synthesis and
survey,'' \textit{European Journal of Operational Research}, vol.~165,
no.~1, pp.~1--19, 2005.

\bibitem{Hakimi1964}
S.~L. Hakimi, ``Optimum locations of switching centers and the absolute
centers and medians of a graph,'' \textit{Operations Research},
vol.~12, no.~3, pp.~450--459, 1964.

\bibitem{Toregas1971}
C.~Toregas, R.~Swain, C.~ReVelle, and L.~Bergman, ``The location of
emergency service facilities,'' \textit{Operations Research}, vol.~19,
no.~6, pp.~1363--1373, 1971.

\bibitem{Church1974}
R.~Church and C.~ReVelle, ``The maximal covering location problem,''
\textit{Papers in Regional Science}, vol.~32, no.~1, pp.~101--118,
1974.

\bibitem{Owen1998}
S.~H. Owen and M.~S. Daskin, ``Strategic facility location: A review,''
\textit{European Journal of Operational Research}, vol.~111, no.~3,
pp.~423--447, 1998.

\bibitem{Drezner2024}
Z.~Drezner and H.~A. Eiselt, ``Competitive location models: A review,''
\textit{European Journal of Operational Research}, vol.~316, no.~1,
pp.~5--18, 2024.

\bibitem{Farahani2012}
R.~Z. Farahani, N.~Asgari, N.~Heidari, M.~Hosseininia, and M.~Goh,
``Covering problems in facility location: A review,''
\textit{Computers \& Industrial Engineering}, vol.~62, no.~2,
pp.~368--407, 2012.

\bibitem{Morrill1977}
R.~L. Morrill and J.~Symons, ``Efficiency and equity aspects of optimum
location,'' \textit{Geographical Analysis}, vol.~9, no.~3,
pp.~215--225, 1977.

\bibitem{Huff1964}
D.~L. Huff, ``Defining and estimating a trading area,''
\textit{Journal of Marketing}, vol.~28, no.~3, pp.~34--38, 1964.

\bibitem{Wilson1971}
A.~G. Wilson, ``A family of spatial interaction models, and associated
developments,'' \textit{Environment and Planning A}, vol.~3, no.~1,
pp.~1--32, 1971.

\bibitem{Nakanishi1974}
M.~Nakanishi and L.~G. Cooper, ``Parameter estimation for a
multiplicative competitive interaction model---Least squares approach,''
\textit{Journal of Marketing Research}, vol.~11, no.~3,
pp.~303--311, 1974.

\bibitem{Fotheringham1983}
A.~S. Fotheringham, ``A new set of spatial-interaction models: The
theory of competing destinations,'' \textit{Environment and Planning A},
vol.~15, no.~1, pp.~15--36, 1983.

\bibitem{Sevtsuk2017}
A.~Sevtsuk and R.~Kalvo, ``Patronage of urban commercial clusters: A
network-based extension of the Huff model for balancing location and
size,'' \textit{City Form Lab Working Paper}, 2017.

\bibitem{Suhara2019}
Y.~Suhara, M.~Bahrami, B.~Bozkaya, and A.~S. Pentland, ``Validating
gravity-based market share models using large-scale transactional
data,'' \textit{arXiv preprint arXiv:1902.03488}, 2019.

\bibitem{Liang2020}
Y.~Liang, S.~Gao, Y.~Cai, N.~Z. Foutz, and L.~Wu, ``Calibrating the
dynamic Huff model for business analysis using location big data,''
\textit{Transactions in GIS}, vol.~24, no.~3, pp.~681--703, 2020.

\bibitem{Verhetsel2022}
A.~Verhetsel, J.~Beckers, and J.~Cant, ``Regional retail landscapes
emerging from spatial network analysis,'' \textit{Regional Studies},
vol.~56, no.~11, pp.~1829--1844, 2022.

\bibitem{Hansen1959}
W.~G. Hansen, ``How accessibility shapes land use,''
\textit{Journal of the American Institute of Planners}, vol.~25,
no.~2, pp.~73--76, 1959.

\bibitem{Geurs2004}
K.~T. Geurs and B.~van Wee, ``Accessibility evaluation of land-use and
transport strategies: Review and research directions,''
\textit{Journal of Transport Geography}, vol.~12, no.~2,
pp.~127--140, 2004.

\bibitem{Luo2003}
W.~Luo and F.~Wang, ``Measures of spatial accessibility to health care
in a GIS environment: Synthesis and a case study in the Chicago
region,'' \textit{Environment and Planning B: Planning and Design},
vol.~30, no.~6, pp.~865--884, 2003.

\bibitem{Wang2012}
F.~Wang, ``Measurement, optimization, and impact of health care
accessibility: A methodological review,'' \textit{Annals of the
Association of American Geographers}, vol.~102, no.~5,
pp.~1104--1112, 2012.

\bibitem{Xu2020}
Y.~Xu, L.~E. Olmos, S.~Abbar, and M.~C. Gonz\'{a}lez,
``Deconstructing laws of accessibility and facility distribution in
cities,'' \textit{Science Advances}, vol.~6, no.~37, p.~eabb4112,
2020.

\bibitem{Kharel2024}
S.~Kharel, S.~Sharifiasl, and Q.~Pan, ``Examining food access equity by
integrating grocery store pricing into spatial accessibility measures,''
\textit{Transportation Research Record}, vol.~2679, no.~1,
pp.~1223--1244, 2024.

\bibitem{Apparicio2007}
P.~Apparicio, M.-S. Cloutier, and R.~Shearmur, ``The case of
Montreal's missing food deserts: Evaluation of accessibility to food
supermarkets,'' \textit{International Journal of Health Geographics},
vol.~6, no.~1, p.~4, 2007.

\bibitem{Beaulac2009}
J.~Beaulac, E.~Kristjansson, and S.~Cummins, ``A systematic review of
food deserts, 1966--2007,'' \textit{Preventing Chronic Disease},
vol.~6, no.~3, p.~A105, 2009.

\bibitem{VerPloeg2012}
M.~Ver Ploeg, V.~Breneman, P.~Dutko, R.~Williams, S.~Snyder,
C.~Dicken, and P.~Kaufman, ``Access to affordable and nutritious food:
Updated estimates of distance to supermarkets using 2010 data,''
\textit{USDA Economic Research Report No.~143}, 2012.

\bibitem{Jin2022}
J.~Jin and Y.~Lu, ``Multi-mode food accessibility: A case study of
Charlotte, NC,'' \textit{Applied Geography}, vol.~145, p.~102749, 2022.

\bibitem{Dijkstra1959}
E.~W. Dijkstra, ``A note on two problems in connexion with graphs,''
\textit{Numerische Mathematik}, vol.~1, no.~1, pp.~269--271, 1959.

\bibitem{Hart1968}
P.~E. Hart, N.~J. Nilsson, and B.~Raphael, ``A formal basis for the
heuristic determination of minimum cost paths,'' \textit{IEEE Trans.\
Systems Science and Cybernetics}, vol.~4, no.~2, pp.~100--107, 1968.

\bibitem{Geisberger2008}
R.~Geisberger, P.~Sanders, D.~Schultes, and D.~Delling, ``Contraction
hierarchies: Faster and simpler hierarchical routing in road networks,''
in \textit{Int.\ Workshop on Experimental and Efficient Algorithms},
Springer, 2008, pp.~319--333.

\bibitem{Delling2011}
D.~Delling, A.~V. Goldberg, I.~Razenshteyn, and R.~F. Werneck,
``PHAST: Hardware-accelerated shortest path trees,'' in \textit{Proc.\
IEEE Int.\ Symp.\ Parallel and Distributed Processing}, 2011,
pp.~921--931.

\bibitem{Bast2016}
H.~Bast, D.~Delling, A.~Goldberg, M.~M\"{u}ller-Hannemann, T.~Pajor,
P.~Sanders, D.~Wagner, and R.~Werneck, ``Route planning in
transportation networks,'' in \textit{Algorithm Engineering: Selected
Results and Surveys}, Springer, 2016, pp.~19--80.

\bibitem{Luxen2011}
D.~Luxen and C.~Vetter, ``Real-time routing with OpenStreetMap data,''
in \textit{Proc.\ 19th ACM SIGSPATIAL Int.\ Conf.\ Advances in
Geographic Information Systems}, 2011, pp.~513--516.

\bibitem{Kang2020}
J.-Y. Kang, A.~Michels, F.~Lyu, S.~Wang, N.~Agbodo, V.~L. Freeman,
and S.~Wang, ``Rapidly measuring spatial accessibility of COVID-19
healthcare resources: A case study of Illinois, USA,''
\textit{International Journal of Health Geographics}, vol.~19, no.~1,
pp.~1--17, 2020.

\bibitem{Horton2025}
P.~Horton, J.~Koh, V.~S. Longman, and A.~Lodi, ``Optimising access
to essential services using the Kolm-Pollak equally distributed
equivalent,'' \textit{Nature Communications}, vol.~16, p.~1234, 2025.

\bibitem{Storn1997}
R.~Storn and K.~Price, ``Differential evolution---A simple and
efficient heuristic for global optimization over continuous spaces,''
\textit{Journal of Global Optimization}, vol.~11, no.~4,
pp.~341--359, 1997.

\bibitem{Nelder1965}
J.~A. Nelder and J.~R. Mead, ``A simplex method for function
minimization,'' \textit{The Computer Journal}, vol.~7, no.~4,
pp.~308--313, 1965.

\bibitem{Virtanen2020}
P.~Virtanen, R.~Gommers, T.~E. Oliphant, M.~Haberland, T.~Reddy,
D.~Cournapeau, E.~Burovski, P.~Peterson, W.~Weckesser, J.~Bright,
S.~J. van~der Walt, M.~Brett, J.~Wilson, K.~J. Millman, N.~Mayorov,
A.~R.~J. Nelson, E.~Jones, R.~Kern, E.~Larson, C.~J. Carey, I.~Polat,
Y.~Feng, E.~W. Moore, J.~VanderPlas, D.~Laxalde, J.~Perktold,
R.~Cimrman, I.~Henriksen, E.~Q. Quintero, C.~M. Harris, A.~H.
Archibald, A.~H. Ribeiro, F.~Pedregosa, and P.~van Mulbregt,
``SciPy 1.0: Fundamental algorithms for scientific computing in
Python,'' \textit{Nature Methods}, vol.~17, pp.~261--272, 2020.

\bibitem{Boeing2017}
G.~Boeing, ``OSMnx: New methods for acquiring, constructing, analyzing,
and visualizing complex street networks,'' \textit{Computers,
Environment and Urban Systems}, vol.~65, pp.~126--139, 2017.

\bibitem{OSM2024}
OpenStreetMap Contributors, ``Planet dump,'' [Online]. Available:
\url{https://planet.openstreetmap.org}, 2024.

\bibitem{Census2023}
U.S. Census Bureau, ``American Community Survey 5-Year Estimates,
2018--2022,'' U.S. Department of Commerce, Washington, DC, 2023.

\end{thebibliography}

\end{document}